\documentclass[a4paper,11pt]{book}%,twocolumn
%\documentclass[a4paper,11pt]{article}
%% packages

\usepackage{blindtext} % needed for creating dummy text passages
%\usepackage{ngerman} % needed for German default language
\usepackage{amsmath} % needed for command eqref
\usepackage{amssymb} % needed for math fonts
\usepackage[colorlinks=true,breaklinks]{hyperref} % needed for creating hyperlinks in the document, the option colorlinks=true gets rid of the awful boxes, breaklinks breaks lonkg links (list of figures), and ngerman sets everything for german as default hyperlinks language
\usepackage[hyphenbreaks]{breakurl} % ben�tigt f�r das Brechen von URLs in Literaturreferenzen, hyphenbreaks auch bei links, die �ber eine Seite gehen (mit hyphenation).
\usepackage{xcolor}
\definecolor{c1}{rgb}{0,0,1} % blue
\definecolor{c2}{rgb}{0,0.3,0.9} % light blue
\definecolor{c3}{rgb}{0.3,0,0.9} % red blue
\hypersetup{
    linkcolor={c1}, % internal links
    citecolor={c2}, % citations
    urlcolor={c3} % external links/urls
}
%\usepackage{cite} % needed for cite
\usepackage[square,authoryear]{natbib} % needed for cite and abbrvnat bibliography style
\usepackage[nottoc]{tocbibind} % needed for displaying bibliography and other in the table of contents
\usepackage{graphicx} % needed for \includegraphics 
\usepackage{longtable} % needed for long tables over pages
\usepackage{bigstrut} % needed for the command \bigstrut
\usepackage{enumerate} % needed for some options in enumerate
%\usepackage{todonotes} % needed for todos
\usepackage{makeidx} % needed for creating an index
\makeindex
\usepackage{gensymb}
\usepackage{url}
\usepackage{psfrag}
\usepackage{multirow}
\usepackage{subfigure}
%% page settings

\usepackage[top=20mm, bottom=20mm,left=20mm,right=20mm]{geometry} % needed for page border settings
\parindent=0mm % for space of first line of new text block
\sloppy % for writing with hyphenless justification (tries to)
\hyphenation{} % use hyphenation of tolerance parametershttp://www.jr-x.de/publikationen/latex/tipps/zeilenumbruch.html
\hyphenpenalty=10000
\exhyphenpenalty=10000
\usepackage{fancyhdr} % needed for head and foot options
%% my macros

%% Text fomats
\newcommand{\tbi}[1]{\textbf{\textit{#1}}}

%% Math fonts
\newcommand{\bbA}{\mathbb{A}}
\newcommand{\bbB}{\mathbb{B}}
\newcommand{\bbC}{\mathbb{C}}
\newcommand{\bbD}{\mathbb{D}}
\newcommand{\bbE}{\mathbb{E}}
\newcommand{\bbF}{\mathbb{F}}
\newcommand{\bbG}{\mathbb{G}}
\newcommand{\bbH}{\mathbb{H}}
\newcommand{\bbI}{\mathbb{I}}
\newcommand{\bbJ}{\mathbb{J}}
\newcommand{\bbK}{\mathbb{K}}
\newcommand{\bbL}{\mathbb{L}}
\newcommand{\bbM}{\mathbb{M}}
\newcommand{\bbN}{\mathbb{N}}
\newcommand{\bbO}{\mathbb{O}}
\newcommand{\bbP}{\mathbb{P}}
\newcommand{\bbQ}{\mathbb{Q}}
\newcommand{\bbR}{\mathbb{R}}
\newcommand{\bbS}{\mathbb{S}}
\newcommand{\bbT}{\mathbb{T}}
\newcommand{\bbU}{\mathbb{U}}
\newcommand{\bbV}{\mathbb{V}}
\newcommand{\bbW}{\mathbb{W}}
\newcommand{\bbX}{\mathbb{X}}
\newcommand{\bbY}{\mathbb{Y}}
\newcommand{\bbZ}{\mathbb{Z}}


\begin{document}
	
\begin{titlepage}
\center % Center everything on the page

%-------------------------------------------------------------------------------------
%	HEADING SECTIONS
%------------------------------------------------------------------------------------
\textbf{\Large L.E. Robotics (Pvt.) Ltd.}\\[0.5cm]
\textbf{\large No. 100/4, Divulapitiya Road, Minuwangoada,	Sri Lanka}\\[2cm]

\includegraphics[width=0.3\textwidth]{figures/logoler}\\[2cm]

	
%-------------------------------------------------------------------------------------
%	TITLE SECTION
%------------------------------------------------------------------------------------
\textbf{\Huge Machine Vision based\\ Real Time Trajectory Generation\\ (MVbR2TG)}\\[6cm]
%\textbf{\Large A comparison}\\[7cm]


%----------------------------------------------------------------------------------------
%	MEMBERS SECTION
%----------------------------------------------------------------------------------------


\vfill

\begin{tabular}[!h]{ l l}
\textbf{\large Prepared by} & {\large Thalagala B.P.}\\
\textbf{\large Last Modified on}&  {\large \today}
\end{tabular}

%\textbf{\large Prepared by}\\[0.5cm]
%{\large Thalagala B.P.}\\[1cm]
%
%%----------------------------------------------------------------------------------------
%%	DATE SECTION
%%----------------------------------------------------------------------------------------
%\textbf{\large Last Modified on}\\[0.5cm]
%\textbf{\Large \today} % Date, change the \today to a set date if you want to be precise

%----------------------------------------------------------------------------------------


\end{titlepage}
\tableofcontents





%\begin{figure}[!h]
%	\centering
%	\includegraphics[scale=0.45]{figures/macllc}
%	\caption{10BASE-T relationship to the ISO/IEC Open Systems Interconnection (OSI) reference model and the IEEE 802.3 CSMA/CD LAN model\cite{main}}
%\end{figure
\chapter{Feasibility Study}
\section{Introduction}
Most of the robotic arms used in industrial environments operate in a pre-programmed cycle. When it comes to the way a human does the same task is much different as the path planning for picking an object may change from cycle to cycle because of the perception obtained through human vision.\\ 

Machine vision is the technology and methods incorporated to mimic the human vision in order to gain the insights about the operating environment of the robotics system. However, when it comes to the real time object detection using machine vision, there is an inevitable trade-off between the accuracy and the speed of the operation. This depends entirely on the used machine vision algorithms and the computational power of the available hardware.

\section{Challenges Encountered when Incorporating  a Vision Unit to a Robotic System}

If the robotic system/ arm in interest is not controlled through a dedicated industrial PC with adequate computational resources, the amount of resources that can be allocated to the vision unit becomes limited.  This will eventually result in great delays (which is not desirable when it comes to real time operations) to produce the required outputs by processing the acquired images through the associated camera.\\

Due to this limitation in computational resources it is not practical to use Deep Learning based algorithms in time critical industrial applications.

\section{Traditional Machine Vision Approach as the Solution}



%\bibliographystyle{plain}
%\bibliography{refer}

%---------------------------------------------------------------------------
\end{document}